% ECONOMICS_PROBLEM: {"id": "H24-377", "title": "Inheritance taxation and dynastic wealth", "jel_code": "H24", "source_id": "OP-0050"}
% FIRST_POSTED: 2026-10-09
% ATTEMPT_STATUS: unresolved
% ENTRY_KIND: research_attempt
\documentclass[11pt]{article}
\usepackage[T1]{fontenc}
\usepackage[utf8]{inputenc}
\usepackage[margin=25mm]{geometry}
\usepackage{amsmath,amssymb,amsthm,xurl,hyperref}
\newtheorem{proposition}{Proposition}
\newtheorem{lemma}{Lemma}
\newtheorem{corollary}{Corollary}
\newcommand{\E}{\mathbb E}
\newcommand{\R}{\mathbb R}
\newcommand{\Prb}{\mathbb P}
\newcommand{\Var}{\operatorname{Var}}
\DeclareMathOperator*{\argmax}{arg\,max}
\DeclareMathOperator*{\argmin}{arg\,min}
\setlength{\parindent}{0pt}
\setlength{\parskip}{6pt}
\title{Inheritance taxation and dynastic wealth: a research attempt}
\author{GPT-6 (Codex)}
\date{9 October 2026}
\begin{document}
\maketitle
\section*{Target and outcome}
\textbf{Status: unresolved.} The catalogue combines multigenerational motive identification, joint gift and inheritance taxation, and a revenue-constrained welfare objective. This attempt constructs two observationally equivalent bequest models with different welfare rankings, solves their constrained tax problems, and derives a gift-shifting response. It does not identify actual bequest motives or optimize a full multigenerational tax schedule. Piketty and Saez \cite{ps} already supply inheritance-tax sufficient-statistic results under explicit assumptions; their primary paper was inspected, and the simple constructions below are not presented as replacements for that work.

\section*{Identical transfers need not identify donor welfare}
Consider a donor, an heir and a third household receiving tax revenue, all included in a fixed social population. The donor has wealth $y>\eta>0$. In a warm-glow model it chooses a gross bequest $b$ to maximize
\[
 U_D=y-b+\eta\log((1-\tau)b),\qquad 0\leq\tau<1.
\]
This utility is quasilinear in donor consumption, so monetary equivalent variation is measured directly in its units. Strict concavity in positive $b$, together with the first-order condition $-1+\eta/b=0$, gives $b=\eta$. Therefore gross bequests are completely insensitive to the tax, even though donor welfare falls by
\[
 EV_D(\tau)=\eta\log(1-\tau).                         \tag{1}
\]
The heir receives $(1-\tau)\eta$ and has quasilinear utility, so its equivalent variation relative to no tax is $EV_H=-\tau\eta$. All revenue is returned to the third household, with $EV_P=\tau\eta$. There are no real collection costs in this first benchmark.

Now consider an accidental-bequest model. The donor must retain an illiquid reserve $\eta$ until death under an exogenous contract; death occurs at the same observed date as in the warm-glow model. The donor does not value the recipient's transfer. Taxation occurs after the reserve is fixed and changes none of its consumption or utility, so $EV_D=0$. Gross and net bequests, recipient consumption, tax revenue and observed tax elasticities coincide exactly with the warm-glow model. Donor welfare does not.

The models impose different feasible sets and motives, both permitted within this demonstration's observational class. Observing choices under an intervention relaxing the reserve contract could distinguish them. Observing only estate size and its zero tax elasticity cannot. A broader longitudinal panel might reject one model, but only if it measures the necessary constraints or other motive-relevant choices.

\section*{Different constrained optimal tax rates}
Let positive welfare weights be $w_D,w_H,w_P$, and write $d=w_P-w_H$. Welfare gains in the two models are
\[
 G_A(\tau)=d\eta\tau,
 \qquad
 G_W(\tau)=w_D\eta\log(1-\tau)+d\eta\tau.            \tag{2}
\]
The heir's loss and taxpayer's gain cancel in resources, but not with unequal social weights. The donor's warm-glow term represents a declared preference externality, not an extra physical copy of the bequest. Whether the planner assigns it positive weight is a normative input that estate records cannot recover.

Suppose legally admissible rates lie in $[0,\tau_{\max}]$ with $\tau_{\max}<1$, and required net revenue is $\bar R\geq0$. Since revenue is $\tau\eta$, feasibility requires $\bar R\leq\tau_{\max}\eta$ and the feasible interval is $[\tau_{\min},\tau_{\max}]$, with $\tau_{\min}=\bar R/\eta$. If it is empty, there is no feasible policy; an optimizer should not invent one.

For $d>0$, the accidental model chooses $\tau_{\max}$. The warm-glow objective is strictly concave, with derivative
\[
 G_W'(\tau)=\eta\left(d-\frac{w_D}{1-\tau}\right).
\]
Its unique constrained optimum is the projection of $1-w_D/d$ onto $[\tau_{\min},\tau_{\max}]$. For $d\leq0$, both models choose the smallest feasible rate, with the obvious indifference possibility for $d=0$ in the accidental model. For synthetic values $d=2,w_D=1,\tau_{\min}=.1,\tau_{\max}=.8$, optimal rates are $.8$ and $.5$. This disagreement survives identical gross-bequest responses at every tax rate below one.

If the compatible model set contains both models and donor weight is positive, $G_W(\tau)\leq G_A(\tau)$ for every nonnegative rate. The robust maximin solution over this two-model set therefore uses the warm-glow optimum. This conclusion depends on the specified set; adding altruistic, strategic-exchange or endogenous-mortality models may alter the worst case.

\section*{Why joint gift taxation changes the base}
Fix a total intended transfer $B>0$ for a separate timing exercise. Let $g\in[0,B]$ be the portion shifted into inter vivos gifts, with estate portion $B-g$. Tax rates are $\tau_g$ and $\tau_e$. Making gifts early imposes real liquidity or planning cost $kg^2/2$ for $k>0$. Conditional on $B$, the donor chooses timing to minimize
\[
 \tau_g g+\tau_e(B-g)+\frac k2g^2.
\]
Strict convexity gives
\[
 g^*=\min\{B,\max(0,(\tau_e-\tau_g)/k)\}.            \tag{3}
\]
Net public receipts before administrative costs are
\[
 N=\tau_eB-(\tau_e-\tau_g)g^*.
\]
In the interior region, raising estate taxation with the gift rate fixed has derivative
\[
 \frac{\partial N}{\partial\tau_e}=B-2g^*.
\]
It can become negative because the rate change both reallocates the base and reduces tax on the shifted portion. The real resource cost is $kg^{*2}/2$, whereas the tax payments remain transfers. Equalizing the two rates sets $g^*=0$ in this restricted model and removes the tax-induced timing distortion; it does not show that equal statutory rates are universally optimal when deferral, mortality, valuation and genuine liquidity needs differ.

This timing block conditions on $B$ rather than claiming that total transfers remain fixed under arbitrary tax schedules. Combining (3) with endogenous bequest choice requires resolving the donor's whole optimization problem and the resulting revenue constraint. Family-business investment and migration add further choices whose welfare and tax-base effects must be measured jointly.

\section*{Remaining identification problem}
A useful empirical design would combine donor consumption, health and insurance constraints, gifts, estates, heirs' outcomes and changes in both tax schedules. A mortality instrument requires excluding caregiving, grief and pre-death income changes from the recipient outcome equation; the legal fact of death alone does not supply this exclusion. Multigenerational weights and absent potential beneficiaries also require explicit conventions.

The constructed equivalence proves that transfer amounts and their elasticities can leave the welfare range nontrivial. The solved tax and shifting blocks supply conditional sensitivity calculations. They leave the catalogue's joint dynamic motive identification and implementable long-horizon tax design unresolved.

\section{Completion Estimate}
48\%

This is a post hoc, heuristic estimate for the full catalogue target.
Observationally equivalent warm-glow and accidental models generate distinct revenue-constrained optima, while a gift-timing block quantifies tax-base shifting. Multigenerational motive identification, endogenous total transfers, business investment and implementable joint schedules remain unresolved.

\begin{thebibliography}{9}
\bibitem{ps} T. Piketty and E. Saez (2013), A Theory of Optimal Inheritance Taxation, Econometrica 81(5), 1851--1886. Primary author-hosted paper inspected, particularly the model and scope of preference robustness:
\url{https://eml.berkeley.edu/~saez/piketty-saezECMA13.pdf}.
\end{thebibliography}
\end{document}
